% =====================================================================
%  style.tex — all formatting lives here. Content files never set
%  spacing, fonts or numbering themselves.
%
%  Every design value is set ONCE in the SETTINGS block below, in Google
%  Docs terms (size, line spacing, space before/after). Everything after
%  it refers back to those settings, so a change there applies
%  everywhere. Below it are only fixed font measurements and standard
%  LaTeX page-break behaviour.
%
%  Needs LuaLaTeX or XeLaTeX (Overleaf: Menu > Compiler > LuaLaTeX).
% =====================================================================

% =====================================================================
%  SETTINGS  (edit here only)
% =====================================================================

% ---- Line spacing (Google Docs: Format > Line spacing) ----
\newcommand{\DocsLineSpacing}{1.15}

% ---- Text sizes, pt ----
\newcommand{\NormalSize}{11}    % Normal text
\newcommand{\TitleSize}{18}     % Title
\newcommand{\SubtitleSize}{9}   % Subtitle: captions, header, page number

% ---- Paragraph spacing, pt (Google Docs "space before/after") ----
\newcommand{\NormalAfter}{12}   % Normal text, after   (before = 0)
\newcommand{\TitleAfter}{4}     % Title, after          (before = 0)
\newcommand{\SubtitleAfter}{12} % Subtitle (captions), after: sub-figure caption → next image
\newcommand{\SubfigureSideGap}{18}% space between side-by-side sub-figures
\newcommand{\CaptionGap}{8}    % space between a figure/table and its own caption
                                % (same for both; edge of image/table → caption text)
% Headings: size, before, after  → applied with \DocsHeading below
\newcommand{\HeadingOne}  {{14}{16}{4}}
\newcommand{\HeadingTwo}  {{12}{4}{4}}
\newcommand{\HeadingThree}{{11}{4}{4}}

% ---- Fonts (files in fonts/) ----
% Weights: Regular, Medium, SemiBold, Bold
\newcommand{\BodyWeight}{Regular}    % normal text
\newcommand{\BoldWeight}{SemiBold}   % headings, caption labels, \textbf
\newcommand{\TitleWeight}{Medium}    % title

% ---- Tables (Google Docs: Table properties) ----
\newcommand{\TableLineSpacing}{1}   % line spacing inside cells (single)
\newcommand{\TableCellPadding}{5}   % cell padding, pt (Docs default 0.07 in)
\newcommand{\TableOuterRule}{1.5}   % top and bottom border, pt, black
\newcommand{\TableHeaderRule}{0.75} % border under the header row, pt, black
\newcommand{\TableRowRule}{0.5}     % border between rows, pt
\newcommand{\TableRowRuleColor}{EFEFEF}  % Docs "light gray 2"
% Automatic width (autotable): the table is as wide as its content
\newcommand{\TableColumnGap}{24}% space between the content of two columns, pt
\newcommand{\TableTextMin}{0.2} % a wrapping text column is never narrower than
                                % this share of the line (or an equal share, if
                                % there are too many); shorter ones stay natural
\newcommand{\TableContinued}{(continued)} % caption text on the next pages of a
                                           % table longer than a page

% ---- Lists and indents (Google Docs: Bullets & numbering > Indentation options)
\newcommand{\ListLabelIndent}{18}  % bullet/number from the margin (0.25 in)
\newcommand{\ListTextIndent}{36}   % text from the margin (0.5 in); each nested
                                   % level and each "Increase indent" adds this

% ---- Alignment of figures, tables and their captions ----
% ieee:     IEEE / NeurIPS / ICML convention (default): narrow figures/tables
%           centred; one-line caption centred, longer caption justified
%           across the full text width
% centered: every caption centred, narrow figures/tables centred
% left:     every caption left-aligned, figures/tables at the left margin
\newcommand{\CaptionStyle}{ieee}

% ---- Table of contents, list of tables/figures (only if used in main.tex)
\newcommand{\TocLineSpacing}{1}    % line spacing of the entries (single)
\newcommand{\TocSectionBefore}{6}  % space before each Heading 1 entry, pt

% ---- Figures/tables ↔ surrounding text, pt (LaTeX only, no Docs equivalent)
\newcommand{\FloatGap}{18}

% ---- Page ----
\newcommand{\PageSetup}{letterpaper, margin=1in,
                        headheight=32pt, headsep=4.6pt, footskip=24pt}

% =====================================================================
%  Below this line: no values, only uses of the settings above
% =====================================================================

% ---------- How Google Docs spaces lines ----------
% Docs line height = line spacing × 1.28 × font size for Lora
% (1.28 = Lora ascender + descender). Gap between paragraphs =
% the LARGER of "after" (paragraph above) and "before" (paragraph below).
\newcommand{\LoraHeight}{1.28}
\newcommand{\LoraDescent}{0.274}
% \DocsPitch[line spacing]{size}; line spacing defaults to \DocsLineSpacing
\NewExpandableDocumentCommand{\DocsPitch}{O{\DocsLineSpacing} m}
  {\fpeval{#1*\LoraHeight*#2}}  % e.g. 1.15 x 1.28 x 11 = 16.19 pt
\newcommand{\DocsBelow}{\fpeval{(\DocsLineSpacing-1)*\LoraHeight/2+\LoraDescent}}
% Set a font size with its Docs line height:  \DocsSize{9}, \DocsSize[1]{11}
\newcommand{\DocsSize}[2][\DocsLineSpacing]{\fontsize{#2}{\DocsPitch[#1]{#2}pt}\selectfont}
% Space between a Docs line's top and its text, per size:  (pitch - size)/2
\newcommand{\DocsTextOffset}[1]{\fpeval{(\DocsPitch{#1}-#1)/2}}

% ---------- How Google Docs places images (measured from your PDF, 0.25 pt steps) ----------
\newcommand{\DocsImageAbove}{3.25}   % paragraph top  → image top
% ---------- How LaTeX places captions (measured, 0.25 pt steps) ----------
% Measured at line spacing 1.5, and at 1.0, 1.15 and 2.0 to get how each
% grows with the caption's line height (\CaptionLineDelta), so the caption
% gaps stay the same at any \DocsLineSpacing.
\newcommand{\CaptionLineDelta}{(\DocsPitch{\SubtitleSize}-\DocsPitch[1.5]{\SubtitleSize})}
\newcommand{\LatexCaptionTop}{\fpeval{5.5+0.69*\CaptionLineDelta}}     % caption skip → caption text top
\newcommand{\LatexCaptionBottom}{\fpeval{3.75+0.30*\CaptionLineDelta}} % caption text bottom → next gap
\newcommand{\LatexTableCaption}{\fpeval{2.75+0.28*\CaptionLineDelta}}  % caption text bottom → table top
\newcommand{\LatexLongTableTop}{\fpeval{2.75+0.17*\CaptionLineDelta}}  % long table caption sits this
                                                                       % much higher than a float's

% ---------- Caption gaps, converted to LaTeX ----------
% image → its caption below  (LaTeX adds \LatexCaptionTop itself)
\newcommand{\FigureCaptionGap}{\fpeval{\CaptionGap-\LatexCaptionTop}}
% sub-figure caption → next image (Docs: Subtitle "after", then the image line)
\newcommand{\SubfigureGap}{\fpeval{\DocsTextOffset{\SubtitleSize}
  +max(\SubtitleAfter,0)+\DocsImageAbove-\LatexCaptionBottom}}
% table caption above → its table  (LaTeX adds \LatexTableCaption itself)
\newcommand{\TableCaptionGap}{\fpeval{\CaptionGap-\LatexTableCaption}}

% ---------- Page ----------
\usepackage{iftex}
\expandafter\usepackage\expandafter[\PageSetup]{geometry}

% ---------- Fonts ----------
\ifPDFTeX
  \errmessage{This template needs LuaLaTeX or XeLaTeX for the Lora fonts.
    On Overleaf: Menu > Compiler > LuaLaTeX, then Recompile}
  \newcommand{\titleweight}{\mdseries}
\else
  \usepackage{lmodern}               % scalable math fonts (no size warnings)
  \usepackage{fontspec}
  \setmainfont{Lora}[
    Path           = fonts/,
    Extension      = .ttf,
    UprightFont    = *-\BodyWeight,
    ItalicFont     = *-Italic,
    BoldFont       = *-\BoldWeight,
    BoldItalicFont = *-\BoldWeight Italic]
  \newfontfamily{\titleweight}{Lora-\TitleWeight}[Path=fonts/, Extension=.ttf,
                                       ItalicFont=Lora-\TitleWeight Italic]
  % Code and URLs: Latin Modern Mono scaled to Lora's lowercase height,
  % so inline code does not look smaller than the text around it
  \setmonofont{lmmono10-regular.otf}[Scale=MatchLowercase,
                                     ItalicFont=lmmono10-italic.otf]
\fi

% ---------- Body text ----------
% Normal text at \NormalSize with the Docs line height.
\makeatletter
\let\Template@normalsize\normalsize
\renewcommand{\normalsize}{\Template@normalsize\DocsSize{\NormalSize}}
\makeatother
\normalsize
\newcommand{\docssmall}{\DocsSize{\SubtitleSize}}
\setlength{\parindent}{0pt}               % no first-line indent
\setlength{\parskip}{\NormalAfter pt}     % Normal "space after"
\AtBeginDocument{\raggedright}            % left-aligned like Google Docs
                                          % (delete this line for justified text)

% ---------- Packages ----------
\usepackage{amsmath}
\usepackage{amssymb}
\usepackage{etoolbox}
\usepackage{graphicx}
\graphicspath{{figures/}}
\usepackage[table]{xcolor}
\usepackage{booktabs}
\usepackage{enumitem}
\usepackage{caption}
\usepackage{subcaption}
\usepackage{titlesec}
\usepackage{fancyhdr}
\usepackage{float}      % [H] placement
\usepackage{flafter}    % a float never appears before its source position
\usepackage{placeins}   % \FloatBarrier
\usepackage{longtable}  % tables longer than a page
\usepackage[hidelinks]{hyperref}
\usepackage[nameinlink, capitalise, noabbrev]{cleveref}   % load last

% ---------- Headings ----------
% Converts Docs size/before/after to LaTeX, assuming Normal text
% above and below the heading.
% \DocsHeading{level}{number format}{\HeadingOne}
\makeatletter
\newcommand{\DocsHeading}[3]{\expandafter\DocsHeadingAux#3{#1}{#2}}
\newcommand{\DocsHeadingAux}[5]{% size, before, after, level, number
  \titleformat{#4}{\DocsSize{#1}\bfseries}{#5}{0.35em}{}[\global\Template@floatlastfalse]%
  \titlespacing*{#4}{0pt}
    {\fpeval{max(\NormalAfter,#2)-\NormalAfter+\DocsBelow*(\NormalSize-#1)}pt}
    {\fpeval{max(#3,0)-\NormalAfter+\DocsBelow*(#1-\NormalSize)}pt}%
}
\DocsHeading{\section}{\thesection.}{\HeadingOne}
\DocsHeading{\subsection}{\thesubsection}{\HeadingTwo}
\DocsHeading{\subsubsection}{\thesubsubsection}{\HeadingThree}
\makeatother

% ---------- Table of contents, lists of tables/figures ----------
% Compact entries: \TocLineSpacing, no paragraph gap, \TocSectionBefore
% above each section entry. Section numbers end with a period, as in the
% headings ("1. Introduction").
\makeatletter
\newcommand{\Template@tocsetup}{%
  \vskip\NormalAfter pt\relax  % the heading's spacing counts on this gap
  \setlength{\parskip}{0pt}%
  \DocsSize[\TocLineSpacing]{\NormalSize}}
\patchcmd{\@starttoc}{\makeatletter}{\makeatletter\Template@tocsetup}{}
  {\errmessage{style.tex: could not patch \string\@starttoc}}
\patchcmd{\l@section}{\addvspace{1.0em \@plus\p@}}{\addvspace{\TocSectionBefore pt}}{}
  {\errmessage{style.tex: could not patch \string\l@section}}
\let\Template@lsection\l@section
\renewcommand{\l@section}[2]{\begingroup
  \let\Template@numberline\numberline
  \renewcommand{\numberline}[1]{\Template@numberline{##1.}}%
  \Template@lsection{#1}{#2}\endgroup}
\makeatother

% ---------- Lists ----------
% Like Google Docs: bullet/number at \ListLabelIndent, text (and wrapped
% lines) at \ListTextIndent; nested levels step in by \ListTextIndent.
% No gap between items; a normal paragraph gap before and after the list.
% beginpenalty: never break a page between a list's intro line and its
% first item (e.g. "...the following:" stays with item 1).
\setlist{itemsep=0pt, topsep=0pt, parsep=0pt, partopsep=0pt, beginpenalty=10000,
         leftmargin=\ListTextIndent pt,
         labelwidth=\fpeval{\ListTextIndent-\ListLabelIndent}pt,
         labelsep=0pt, align=left}
\setlist[itemize,1]{label=\textbullet}
\setlist[itemize,2]{label=$\circ$}
\setlist[itemize,3]{label=$\scriptscriptstyle\blacksquare$}
\setlist[enumerate,1]{label=\arabic*.}
\setlist[enumerate,2]{label=\alph*.}
\setlist[enumerate,3]{label=\roman*.}
\setlist[description]{font=\bfseries, labelwidth=0pt, labelsep=0.5em}

% ---------- Indented paragraphs ----------
% Like "Increase indent" in Google Docs:
%   \begin{indented} ... \end{indented}       one step  (\ListTextIndent)
%   \begin{indented}[2] ... \end{indented}    two steps
\newenvironment{indented}[1][1]
  {\par\setlength{\leftskip}{\fpeval{#1*\ListTextIndent}pt}}
  {\par}
\renewenvironment{quote}{\begin{indented}}{\end{indented}}

% ---------- Equations ----------
% Extra space between lines of align/gather so stacked fractions don't touch
\setlength{\jot}{\fpeval{\NormalAfter/2}pt}

% ---------- Footnotes ----------
\renewcommand{\footnotesize}{\docssmall}

% ---------- Captions ----------
\makeatletter
% Small text, bold label: "Table 1." / "Figure 1a."
\DeclareCaptionFont{docssmall}{\docssmall}
% Alignment from \CaptionStyle: caption text, and narrow figures/tables
\ifdefstring{\CaptionStyle}{left}
  {\newcommand{\Template@capalign}{justification=raggedright, singlelinecheck=false}%
   \newcommand{\Template@floatalign}{\raggedright}}
  {\ifdefstring{\CaptionStyle}{ieee}
    {\newcommand{\Template@capalign}{justification=justified, singlelinecheck=true}}
    {\newcommand{\Template@capalign}{justification=centering, singlelinecheck=false}}%
   \newcommand{\Template@floatalign}{\centering}}
% One caption style, defined once and used by captions and sub-captions
% [ ] = how a caption that fits on one line is set (used with singlelinecheck)
\edef\Template@tmp{\noexpand\DeclareCaptionStyle{docs}[justification=centering]{font=docssmall,
  labelfont=bf, labelsep=period, skip=\FigureCaptionGap pt, \Template@capalign}}
\Template@tmp
\captionsetup{style=docs}
% Tables: caption ABOVE.  Figures: caption BELOW.
\captionsetup[table]{position=top, skip=\TableCaptionGap pt}
\captionsetup[figure]{position=bottom}
% Sub-figures: the same style, label "Figure 1a"
\DeclareCaptionLabelFormat{subfull}{Figure~\thefigure#2}
\captionsetup[sub]{style=docs, labelformat=subfull}
\renewcommand{\thesubfigure}{\alph{subfigure}}
\crefname{subfigure}{Figure}{Figures}
% Between stacked sub-figures:  \end{subfigure} \subfiguregap \begin{subfigure}
\newcommand{\subfiguregap}{\par\vspace{\SubfigureGap pt}}
% Width of one of n side-by-side sub-figures:  \begin{subfigure}[t]{\subwidth{2}}
% (n widths + (n-1) \SubfigureSideGap = the line); separate them with \hfill
\newcommand{\subwidth}[1]{\dimexpr(\linewidth-\SubfigureSideGap pt*\numexpr#1-1\relax)/#1\relax}
\makeatother

% ---------- Narrow figures and tables ----------
% Centred or at the left margin, from \CaptionStyle (full-width ones are
% unaffected). LaTeX resets alignment inside the float box, so it is set there.
\makeatletter
\g@addto@macro\@floatboxreset{\Template@floatalign}
\makeatother

% ---------- Figure/table placement ----------
% Standard LaTeX placement with two safeguards:
%  - a figure/table goes where it is written if it fits ("here"), otherwise
%    to the top of the next page, and the text fills the gap (no empty half
%    pages);
%  - it never appears before the place it is written (flafter), and never
%    moves past the next numbered section, e.g. "3." (\FloatBarrier).
% To force one to stay exactly where it is written: \begin{figure}[H].
\setlength{\floatsep}{\FloatGap pt plus 2pt minus 2pt}
\setlength{\textfloatsep}{\floatsep}
\setlength{\intextsep}{\floatsep}
% Pages holding only floats: stack from the top instead of spreading out.
\makeatletter
\setlength{\@fptop}{0pt}
\setlength{\@fpsep}{\floatsep}
\setlength{\@fpbot}{0pt plus 1fil}
\makeatother
% Space around a figure/table placed in the text: \FloatGap above and below,
% measured like Google Docs from the edge of the text line (baseline +
% \Template@linebelow), not from the baseline. Without this the gap above
% is about 6 pt smaller than the gap below, and a heading after a figure
% gets an extra paragraph gap.
\makeatletter
\newdimen\Template@linebelow
\setlength{\Template@linebelow}{\fpeval{\DocsBelow*\NormalSize}pt}
\newif\ifTemplate@afterheading
\newif\ifTemplate@floatlast   % last thing placed was a figure/table, not a heading
% "here" placement (LaTeX's own)
\patchcmd{\@addtocurcol}
  {\vskip \intextsep \box\@currbox \penalty\interlinepenalty \vskip\intextsep
   \ifnum\outputpenalty <-\@Mii \vskip -\parskip\fi}
  {\vskip \intextsep \vskip\dimexpr\Template@linebelow-\@pagedp\relax
   \box\@currbox \penalty\interlinepenalty
   \vskip\intextsep \vskip\dimexpr\@pagedp-\Template@linebelow\relax
   \ifnum\outputpenalty <-\@Mii \kern-\parskip\fi}{}
  {\errmessage{style.tex: could not patch \string\@addtocurcol}}
% [H] placement (float package); after a heading the heading's own spacing
% already applies, so only the space below is adjusted
\patchcmd{\float@endH}{\@endfloatbox\vskip\intextsep}
  {\@endfloatbox\vskip\intextsep
   \ifTemplate@afterheading\else\ifdim\prevdepth>-1000pt
     \vskip\dimexpr\Template@linebelow-\prevdepth\relax\fi\fi}{}
  {\errmessage{style.tex: could not patch \string\float@endH}}
\apptocmd{\float@endH}{\prevdepth\Template@linebelow\kern-\parskip
  \global\Template@floatlasttrue}{}
  {\errmessage{style.tex: could not patch \string\float@endH}}
\makeatother
\floatplacement{figure}{htbp}
\floatplacement{table}{htbp}
\renewcommand{\topfraction}{0.85}       % a page top may be mostly figure
\renewcommand{\bottomfraction}{0.6}
\renewcommand{\textfraction}{0.15}      % keep at least 15% text on a page
\renewcommand{\floatpagefraction}{0.75} % figure-only pages must be fairly full
\setcounter{topnumber}{3}
\setcounter{bottomnumber}{2}
\setcounter{totalnumber}{4}
% A figure/table written directly after a heading (no text in between) stays
% attached to that heading: it is placed right there, and if both don't fit,
% they move to the next page together.
\makeatletter
\AtBeginEnvironment{figure}{\if@nobreak\def\fps@figure{H}\global\Template@afterheadingtrue
  \else\global\Template@afterheadingfalse\fi}
\AtBeginEnvironment{table}{\if@nobreak\def\fps@table{H}\global\Template@afterheadingtrue
  \else\global\Template@afterheadingfalse\fi}
\makeatother

% ---------- Page breaks ----------
% Never leave a single line of a paragraph alone at the bottom of a page
% (orphan) or at the top of the next one (widow), like Google Docs'
% "widow/orphan control". Pages may end slightly early instead.
\clubpenalty=10000
\widowpenalty=10000
\displaywidowpenalty=10000
\raggedbottom
% Headings are never left alone at the bottom of a page ("keep with next"):
% a heading needs room for itself plus about two lines, otherwise new page.
% A heading directly after another heading skips the check (it would break
% the page between the two); the headings above it already keep the block
% together.
\usepackage{needspace}
\makeatletter
\newcommand{\Template@keepwithnext}[1]{%
  \if@nobreak\ifTemplate@floatlast\Needspace{#1\baselineskip}\fi
  \else\Needspace{#1\baselineskip}\fi}
\let\Template@section\section
\renewcommand{\section}{\FloatBarrier\Template@keepwithnext{3}\Template@section}
\let\Template@subsection\subsection
\renewcommand{\subsection}{\Template@keepwithnext{3}\Template@subsection}
\let\Template@subsubsection\subsubsection
\renewcommand{\subsubsection}{\Template@keepwithnext{2.5}\Template@subsubsection}
\makeatother

% ---------- Tables ----------
% Docs cell = padding + one line (at \TableLineSpacing) + padding, with the
% text centred; borders sit between rows. Each row gets exactly that height.
\newcommand{\CellEdge}{\fpeval{\TableCellPadding
  +(\TableLineSpacing-1)*\LoraHeight*\NormalSize/2}}  % padding + half the extra line space
\newcommand{\CellAbove}{\fpeval{\CellEdge+(\LoraHeight-\LoraDescent)*\NormalSize}}
\newcommand{\CellBelow}{\fpeval{\CellEdge+\LoraDescent*\NormalSize}}
\makeatletter
\newcommand{\Template@tablestrut}{%
  % cell text at table line spacing (also for text that wraps in a cell)
  \DocsSize[\TableLineSpacing]{\NormalSize}%
  \renewcommand{\arraystretch}{1}\setlength{\extrarowheight}{0pt}%
  \setbox\strutbox\hbox{\vrule height \CellAbove pt depth \CellBelow pt width 0pt}}
\makeatother
\AtBeginEnvironment{tabular}{\csname Template@tablestrut\endcsname}
\AtBeginEnvironment{tabular*}{\csname Template@tablestrut\endcsname}
\AtBeginEnvironment{tabulary}{\csname Template@tablestrut\endcsname}
% Between the content of two columns \TableColumnGap; at the table's left
% and right edges the cell padding (autotable adds it with \TemplatePad)
\setlength{\tabcolsep}{\fpeval{\TableColumnGap/2}pt}
\newcommand{\TemplatePad}{\hspace{\TableCellPadding pt}}
% Borders: no extra space around rules, padding is already in the rows
\setlength{\aboverulesep}{0pt}
\setlength{\belowrulesep}{0pt}
\setlength{\heavyrulewidth}{\TableOuterRule pt}   % \toprule, \bottomrule
\setlength{\lightrulewidth}{\TableHeaderRule pt}  % \midrule
\definecolor{TableRowRuleColor}{HTML}{\TableRowRuleColor}
% Light line between body rows:  Row 1 \\ \rowline  Row 2 \\
\newcommand{\rowline}{\arrayrulecolor{TableRowRuleColor}%
  \midrule[\TableRowRule pt]\arrayrulecolor{black}}

% ---------- Automatic table width ----------
%   \begin{autotable}{L R R}  ...rows...  \end{autotable}
% Column types:  L = text (left), R = numbers (right), C = centred.
% The table is as wide as its content, with \TableColumnGap between the
% content of neighbouring columns. If that is wider than the line, it is set
% to full width and its text (L) columns wrap, each getting width in
% proportion to its text; R/C columns keep their natural width.
% Tables narrower than the line are centred (see Narrow figures and tables).
\usepackage{environ}
\usepackage{tabulary}
\setlength{\tymin}{\fpeval{\TableTextMin}\linewidth}
% With too many text columns for that share, tabulary gives them equal widths;
% it warns each time, which is expected here (e.g. four text columns on A4)
\makeatletter
\patchcmd{\endtabulary}{\TY@warn{tymin too large (\the\tymin), resetting to \the\dimen@}}{}{}
  {\errmessage{style.tex: could not patch \string\endtabulary}}
\makeatother
\ExplSyntaxOn
\tl_new:N \l_template_plainspec_tl   % L R C → l r c   (natural width)
\tl_new:N \l_template_wrapspec_tl    % L stays L (wraps), R C → r c
\int_new:N \l_template_ntext_int
\cs_new_protected:Npn \TemplateSpecs #1
  {
    \tl_set:Nn \l_template_wrapspec_tl {#1}
    \tl_replace_all:Nnn \l_template_wrapspec_tl {R} {r}
    \tl_replace_all:Nnn \l_template_wrapspec_tl {C} {c}
    \tl_set_eq:NN \l_template_plainspec_tl \l_template_wrapspec_tl
    \tl_replace_all:Nnn \l_template_plainspec_tl {L} {l}
    \regex_count:nnN {L} {#1} \l_template_ntext_int
  }
\cs_new:Npn \TemplateNtext { \int_use:N \l_template_ntext_int }
\cs_new_protected:Npn \TemplatePlainTabular #1
  {
    \exp_args:Nnx \begin {tabular}
      { @{\exp_not:N\TemplatePad} \l_template_plainspec_tl @{\exp_not:N\TemplatePad} }
    #1 \end{tabular}
  }
\cs_new_protected:Npn \TemplateStretchedTabular #1 #2
  {
    \exp_args:Nnnx \begin {tabular*} {#1}
      { @{\exp_not:N\extracolsep{\exp_not:N\fill}\exp_not:N\TemplatePad}
        \l_template_plainspec_tl @{\exp_not:N\TemplatePad} }
    #2 \end{tabular*}
  }
\cs_new_protected:Npn \TemplateWrappedTabular #1
  {
    \exp_args:Nnnx \begin {tabulary} {\linewidth}
      { @{\exp_not:N\TemplatePad} \l_template_wrapspec_tl @{\exp_not:N\TemplatePad} }
    #1 \end{tabulary}
  }
\ExplSyntaxOff
\makeatletter
\NewEnviron{autotable}[1]{%
  \TemplateSpecs{#1}%
  \setbox\@ne\hbox{\TemplatePlainTabular{\BODY}}%
  \ifTemplate@long
    \expandafter\TemplateLongTabular\expandafter{\BODY}% longer than a page (see below)
  \else\ifdim\wd\@ne>\linewidth
    \ifnum\TemplateNtext>0
      \TemplateWrappedTabular{\BODY}%       too long: wrap text columns
    \else
      \TemplateStretchedTabular{\linewidth}{\BODY}% numbers only: cannot wrap
    \fi
  \else
    \usebox\@ne%                            natural width
  \fi\fi
}
\makeatother

% ---------- Tables longer than a page ----------
% A table whose caption and rows do not fit on one page is not a float: it
% is set where it is written and continues on the next pages, each part
% under "Table N. (continued)" (\TableContinued) with the header rows
% repeated. It is written like any other table; the switch is automatic.
% Column widths are the same as for a short table (natural, stretched, or
% the widths tabulary chooses). A page never ends with only the header and
% one row, or starts with only the last row.
\makeatletter
\newif\ifTemplate@long
\newsavebox\Template@measurebox
% Every table first collects its content and measures it, caption included
\AtBeginEnvironment{table}{\let\table\Template@table}
\def\Template@table{\Collect@Body\Template@tablebody}
\long\def\Template@tablebody#1{\Template@tableopt#1\@nil}
\def\Template@tableopt{\@ifnextchar[\Template@tableoptA{\Template@tableoptA[]}}
\long\def\Template@tableoptA[#1]#2\@nil{%
  \edef\Template@count{\the\c@table}%
  \setbox\Template@measurebox\vbox{\hsize\columnwidth\linewidth\hsize
    \@parboxrestore\@floatboxreset\def\@captype{table}\let\label\@gobble
    \captionsetup{hypcap=false}#2}%
  \global\c@table\Template@count\relax
  \ifdim\dimexpr\ht\Template@measurebox+\dp\Template@measurebox\relax>\textheight
    \expandafter\@firstoftwo\else\expandafter\@secondoftwo\fi
  {\par\Template@longtrue\let\endtable\relax
   \global\let\Template@cap\@empty\global\let\Template@label\@empty
   \let\Template@reallabel\label
   \def\caption{\@dblarg\Template@storecaption}%
   \def\label##1{\gdef\Template@label{##1}}%
   \captionsetup{hypcap=false}% the link goes to the table start instead
   #2}%
  {\ifblank{#1}{\@float{table}#2}{\@float{table}[#1]#2}}}% a normal float
\long\def\Template@storecaption[#1]#2{\gdef\Template@lotcap{#1}\gdef\Template@cap{#2}}
% Column widths of a wrapping table: run tabulary's measuring pass and keep
% the widths it computes (\TY@F1, \TY@F2 ...) instead of setting the table
\newenvironment{Template@tywidths}
  {\let\TY@final\Template@tycapture\let\endTY@final\relax\TY@tabular}
  {\endtabulary}
\long\def\Template@tycapture#1#2\endTY@final{%
  \count@\z@
  \loop\ifnum\count@<\Template@ncols\relax
    \advance\count@\@ne
    \expandafter\xdef\csname Template@w\the\count@\endcsname
      {\the\dimexpr\csname TY@F\the\count@\endcsname\relax}%
  \repeat}
% Spacing like a figure/table placed in the text (\float@endH above). After
% a heading: no page break between the two; if the header and first rows do
% not fit, both go to the next page (longtable would break after the heading)
\patchcmd{\LT@start}{\endgraf\penalty\z@\vskip\LTpre}
  {\endgraf\Template@ltbreak\vskip\LTpre}{}
  {\errmessage{style.tex: could not patch \string\LT@start}}
\patchcmd{\LT@start}{\ifdim \dimen@>\z@ \vfil\break}
  {\ifdim \dimen@>\z@ \Template@ltfull}{}
  {\errmessage{style.tex: could not patch \string\LT@start}}
\def\Template@ltbreak{\penalty\z@}
\def\Template@ltfull{\vfil\break}
\newcommand{\Template@ltspacing}{%
  % fixed (no shrink): a long table fills its page, so it would use the shrink
  \setlength{\LTpre}{\dimexpr\FloatGap pt+\LatexLongTableTop pt\relax}%
  \setlength{\LTpost}{\intextsep}%
  \ifTemplate@afterheading   % placed like [H], which adds \lineskip above
    \addtolength{\LTpre}{\lineskip}%
    \def\Template@ltbreak{\nobreak}\let\Template@ltfull\relax
  \else\ifdim\prevdepth>-1000pt
    \addtolength{\LTpre}{\dimexpr\Template@linebelow-\prevdepth\relax}\fi\fi}
\newcommand{\Template@ltafter}{%
  \prevdepth\Template@linebelow\kern-\parskip\global\Template@floatlasttrue}
\AtBeginEnvironment{longtable}{\Template@tablestrut}
% \rowline in a long table: the same light line, drawn as a row across all
% columns like booktabs does in longtable, with the only page break point
% between two rows right after it (none in the "nb" version)
\newcommand{\Template@ltrule}[1]{\noalign{\nobreak}\multispan\LT@cols
  {\color{TableRowRuleColor}\leaders\hrule\@height\TableRowRule pt\hfill\kern\z@}%
  \cr\noalign{#1}}
\newcommand{\Template@ltrowline}{\Template@ltrule{\penalty\z@}}
\newcommand{\Template@ltrowlinenb}{\Template@ltrule{\nobreak}}
% A caption row gets the caption's own strut instead of the cell strut, so
% it starts where a figure/table caption would (with the cell strut it sits
% too low; with none, its text reaches above the row, out of the text area)
\newsavebox\Template@arstrut
\newcommand{\Template@nostrut}{\noalign{\global\setbox\Template@arstrut\copy\@arstrutbox
  \begingroup\docssmall
  \global\setbox\@arstrutbox\hbox{\vrule\@height\ht\strutbox\@width\z@}\endgroup}}
\newcommand{\Template@strut}{\noalign{\global\setbox\@arstrutbox\box\Template@arstrut}}
\ExplSyntaxOn
\tl_new:N \l_template_head_tl
\tl_new:N \l_template_rows_tl
\tl_new:N \l_template_ltspec_tl
\tl_new:N \l_template_lt_tl
\seq_new:N \l_template_rows_seq
\int_new:N \l_template_nrows_int
\cs_new_protected:Npn \__template_split_head:w #1 \midrule #2 \q_stop
  { \tl_set:Nn \l_template_head_tl {#1} \tl_set:Nn \l_template_rows_tl {#2} }
\cs_new_protected:Npn \__template_strip_bottom:w #1 \bottomrule #2 \q_stop
  { \tl_set:Nn \l_template_rows_tl {#1} }
\cs_new_protected:Npn \TemplateLongTabular #1
  {
    % column spec and horizontal position, as autotable would choose them
    \dim_compare:nNnTF { \box_wd:N 1 } > { \linewidth }
      {
        \int_compare:nNnTF { \TemplateNtext } > { 0 }
          {
            \cs_set:Npx \Template@ncols { \tl_count:N \l_template_wrapspec_tl }
            \hbox_set:Nn \l_tmpa_box
              {
                \Template@tablestrut
                \exp_args:Nnnx \begin {Template@tywidths} {\linewidth}
                  { @{\exp_not:N\TemplatePad} \l_template_wrapspec_tl @{\exp_not:N\TemplatePad} }
                #1 \end{Template@tywidths}
              }
            \tl_set:Nn \l_template_ltspec_tl { @{\TemplatePad} }
            \int_zero:N \l_tmpa_int
            \tl_map_inline:Nn \l_template_wrapspec_tl
              {
                \int_incr:N \l_tmpa_int
                \str_if_eq:nnTF {##1} {L}
                  {
                    \tl_put_right:Nx \l_template_ltspec_tl
                      { > { \exp_not:N \raggedright \exp_not:N \arraybackslash }
                        p { \use:c { Template@w \int_use:N \l_tmpa_int } } }
                  }
                  { \tl_put_right:Nn \l_template_ltspec_tl {##1} }
              }
            \tl_put_right:Nn \l_template_ltspec_tl { @{\TemplatePad} }
            \setlength{\LTleft}{\fill} \setlength{\LTright}{\fill}
          }
          {
            \tl_set:Nx \l_template_ltspec_tl
              { @{ \exp_not:N \extracolsep { \exp_not:N \fill } \exp_not:N \TemplatePad }
                \l_template_plainspec_tl @{ \exp_not:N \TemplatePad } }
            \setlength{\LTleft}{0pt} \setlength{\LTright}{0pt}
          }
      }
      {
        \tl_set:Nx \l_template_ltspec_tl
          { @{ \exp_not:N \TemplatePad } \l_template_plainspec_tl @{ \exp_not:N \TemplatePad } }
        \str_if_eq:VnTF \CaptionStyle {left}
          { \setlength{\LTleft}{0pt} } { \setlength{\LTleft}{\fill} }
        \setlength{\LTright}{\fill}
      }
    % header rows (up to the first \midrule), body rows, no final \bottomrule
    \tl_if_in:nnTF {#1} { \midrule }
      { \__template_split_head:w #1 \q_stop }
      { \tl_clear:N \l_template_head_tl \tl_set:Nn \l_template_rows_tl {#1} }
    \tl_if_in:NnT \l_template_rows_tl { \bottomrule }
      { \exp_last_unbraced:NV \__template_strip_bottom:w \l_template_rows_tl \q_stop }
    \seq_set_split:NnV \l_template_rows_seq { \\ } \l_template_rows_tl
    \seq_pop_right:NN \l_template_rows_seq \l_tmpa_tl
    \tl_if_blank:VF \l_tmpa_tl { \seq_put_right:NV \l_template_rows_seq \l_tmpa_tl }
    \int_set:Nn \l_template_nrows_int { \seq_count:N \l_template_rows_seq }
    % caption + header on the first page, "(continued)" + header on the next
    \tl_set:Nx \l_template_lt_tl
      { \exp_not:N \begin {longtable} { \exp_not:V \l_template_ltspec_tl } }
    \tl_if_empty:NF \Template@cap
      {
        \tl_put_right:Nx \l_template_lt_tl
          {
            \exp_not:N \Template@nostrut
            \exp_not:N \caption [ \exp_not:V \Template@lotcap ] { \exp_not:V \Template@cap }
            \tl_if_empty:NF \Template@label
              { \exp_not:N \Template@reallabel { \exp_not:V \Template@label } }
            \exp_not:N \\ \exp_not:N \Template@strut
          }
      }
    \tl_if_empty:NTF \l_template_head_tl
      { \tl_put_right:Nn \l_template_lt_tl { \endfirsthead } }
      { \tl_put_right:Nx \l_template_lt_tl
          { \exp_not:V \l_template_head_tl \exp_not:N \midrule \exp_not:N \endfirsthead } }
    \tl_if_empty:NF \Template@cap
      {
        \tl_put_right:Nn \l_template_lt_tl
          { \Template@nostrut \caption [] { \TableContinued } \\ \Template@strut }
      }
    \tl_if_empty:NTF \l_template_head_tl
      { \tl_put_right:Nn \l_template_lt_tl { \toprule \endhead } }
      { \tl_put_right:Nx \l_template_lt_tl
          { \exp_not:V \l_template_head_tl \exp_not:N \midrule \exp_not:N \endhead } }
    \tl_put_right:Nn \l_template_lt_tl
      { \bottomrule \endfoot \bottomrule \endlastfoot }
    % rows; no page break after the first row or before the last one
    \seq_map_indexed_inline:Nn \l_template_rows_seq
      {
        \bool_lazy_and:nnTF
          {
            \bool_lazy_or_p:nn
              { \int_compare_p:nNn {##1} = {2} }
              { \int_compare_p:nNn {##1} = { \l_template_nrows_int } }
          }
          { \tl_if_head_eq_meaning_p:nN {##2} \rowline }
          {
            \tl_put_right:Nn \l_template_lt_tl { \Template@ltrowlinenb }
            \tl_put_right:Nx \l_template_lt_tl { \exp_not:o { \tl_tail:n {##2} } }
          }
          { \tl_put_right:Nn \l_template_lt_tl {##2} }
        \bool_lazy_or:nnTF
          { \int_compare_p:nNn {##1} = {1} }
          { \int_compare_p:nNn {##1} = { \l_template_nrows_int - 1 } }
          { \tl_put_right:Nn \l_template_lt_tl { \\* } }
          { \tl_put_right:Nn \l_template_lt_tl { \\ } }
      }
    \tl_put_right:Nn \l_template_lt_tl { \end{longtable} }
    \Template@ltspacing
    \cs_set_eq:NN \rowline \Template@ltrowline
    \tl_use:N \l_template_lt_tl
    \Template@ltafter
  }
\cs_generate_variant:Nn \seq_set_split:Nnn { NnV }
\cs_generate_variant:Nn \str_if_eq:nnTF { Vn }
\ExplSyntaxOff
\makeatother

% ---------- Header / footer ----------
\makeatletter
% Page number bottom right. Course + author top right on page 1.
\renewcommand{\headrulewidth}{0pt}
\newcommand{\Template@pagenumber}{\fancyfoot[R]{\docssmall\thepage}}
\fancypagestyle{plain}{\fancyhf{}\Template@pagenumber}
\pagestyle{plain}
\makeatother

% ---------- Title block ----------
\makeatletter
\newcommand{\course}[1]{\def\@course{#1}}
\def\@course{}
\fancypagestyle{firstpage}{%
  \fancyhf{}%
  \fancyhead[R]{\docssmall\@course\\\@author}%
  \Template@pagenumber}
\renewcommand{\maketitle}{%
  \thispagestyle{firstpage}%
  {\centering\DocsSize{\TitleSize}\titleweight \@title\par}%
  % next block's own "before" handles the gap; correct for title size
  \vspace{\fpeval{max(\TitleAfter,\NormalAfter)-\NormalAfter
                  +\DocsBelow*(\TitleSize-\NormalSize)}pt}%
}
\makeatother
